\documentclass{article}
\usepackage{amsmath}
\usepackage{amssymb}
\usepackage{amsthm}
\usepackage{physics}


\newtheorem*{conjecture}{Conjecture}
\newtheorem*{theorem}{Theorem}
\newtheorem*{lemma}{Lemma}


\makeatletter
\newcommand*{\rom}[1]{\expandafter\@slowromancap\romannumeral #1@}
\makeatother


\begin{document}

\section*{1}

\subsection*{(\rom{1})}

We will use the properties of numbers stated in chapter 1:

\begin{equation*}
	ax = a
\end{equation*}
\begin{equation*}
	a^{-1}(ax) = a^{-1}a \ \ \ \mbox{\emph{Multiplying by the inverse}}	
\end{equation*}
\begin{equation*}
	(a^{-1}a)x = 1  \ \ \ \mbox{\emph{Associative law for multiplication}}
\end{equation*}
\begin{equation*}
	1 \cdot x = 1  \ \ \ \mbox{\emph{Existence of multiplicative inverses}}
\end{equation*}
\begin{equation*}
	x = 1  \ \ \ \mbox{\emph{Existence of a multiplicative identity}}
\end{equation*}

and the proof is done.


\subsection*{(\rom{2})}

\begin{equation*}
	(x-y)(x+y) =
\end{equation*}
\begin{equation*}
	=(x-y) \cdot x + (x-y) \cdot y \ \ \ \mbox{\emph{Distributive law}}
\end{equation*}
\begin{equation*}
	=x \cdot x + (-y \cdot x) + x \cdot y + (-y \cdot y) \ \ \
	\mbox{\emph{Distributive law}}
\end{equation*}
\begin{equation*}
	=x \cdot x + (-y \cdot x) + y \cdot x + (-y \cdot y)  \ \ \
	\mbox{\emph{Commutative law for multiplication}}
\end{equation*}
\begin{equation*}
	=x \cdot x +  (-y \cdot y) \ \ \
	\mbox{\emph{Existence of additive inverses}}
\end{equation*}
\begin{equation*}
	= x^2 - y^2
\end{equation*}


\subsection*{(\rom{3})}
We start by subtracting \(y^2\) from both sides of our equation: 
\begin{equation*}
	x^2 = y^2  \ \ \ / +(-y^2) 
\end{equation*}
We will use the property of  existence of additive inverse
\begin{equation*}
	x^2 - y^2 = y^2 - y^2 = 0
\end{equation*}
Now, let us use the property we proved in exercise 1.(II):
\begin{equation*}
	x^2 - y^2 = (x-y)(x+y) = 0
\end{equation*}
From the above, either \((x-y)=0\) or \((x+y) = 0\) so using simple algebra we
end up with
\begin{equation*}
	x = y
\end{equation*}
\begin{equation*}
	x = -y
\end{equation*}
and our proof is done.



\subsection*{(\rom{4})}
Let us use the power of algebra and let use expand the expression on the right
side of equation:
\begin{equation*}
	(x-y)(x^2 + xy + y^2) = x^{3} + x^{2}y + xy^{2} - yx^{2} - xy^{2} -
	y^{3}
\end{equation*}
Now we will apply the properties of real numbers:
\begin{equation*}
	x^{3} + x^{2}y + xy^{2} - x^{2}y - xy^{2} -
	y^{3} \ \ \  \mbox{\emph{Commutative law for multiplication}}
\end{equation*}
\begin{equation*}	
	=x^{3} + x^{2}y - x^{2}y + xy^{2} - xy^{2} -
	y^{3} \ \ \  \mbox{\emph{Commutative law for addition}}
\end{equation*}
\begin{equation*}	
	=x^{3} + 0 + 0 -
	y^{3} \ \ \  \mbox{\emph{Existence of additive inverses}}
\end{equation*}
\begin{equation*}	
	=x^{3} - y^{3} \ \ \  \mbox{\emph{Existence of an additive identity}}
\end{equation*}
and the last step concludes our proof.

\subsection*{(\rom{5})}
We will distribute the expression over (\(x-y)\):
\begin{equation*}
	(x-y) \cdot (x^{n-1} + x^{n-2}y + \cdots + xy^{n-2} + y^{n-1}) = 
\end{equation*}
\begin{multline*}
	= x \cdot (x^{n-1} + x^{n-2}y + \cdots + xy^{n-2} + y^{n-1})\\
	+ \Big(-y \cdot (x^{n-1} + x^{n-2}y + \cdots + xy^{n-2} + y^{n-1}) \Big) = 
\end{multline*}
\begin{multline*}
	= (x \cdot x^{n-1} + x \cdot x^{n-2}y + \cdots + x \cdot y^{n-2} + x \cdot y^{n-1})\\
	+ \Big(-y \cdot x^{n-1} -y \cdot x^{n-2}y - \cdots -y \cdot xy^{n-2} -y \cdot y^{n-1}\Big) = 
\end{multline*}
\begin{multline*}
	= (x^{n} + x^{n-1}y + \cdots + x^{2} \cdot y^{n-2} + x \cdot y^{n-1})\\
	+ \Big(-y \cdot x^{n-1} - x^{n-2}y^{2} - \cdots - xy^{n-1} - y^{n}\Big)
\end{multline*}
Let us use commutative laws (addition and multiplication) to rearrange the terms, and present the equation to find a pattern:
\begin{multline*}
	= (x^{n} + x^{n-1}y + x^{n-2}y^2 + x^{n-3}y^3 + \cdots + x^{2}  y^{n-2} + x  y^{n-1})\\
	+ \Big(-y  x^{n-1} - x^{n-2}y^{2} - x^{n-3}y^{3} - \cdots - x^{2}y^{n-2} - xy^{n-1} - y^{n}\Big) = 
\end{multline*}
\begin{multline*}
	= (x^{n} + x^{n-1}y + x^{n-2}y^2 + x^{n-3}y^3 + \cdots + x^{2}  y^{n-2} + x  y^{n-1})\\
	+ \Big(-y^{n} - x^{n-1}y - x^{n-2}y^{2} - x^{n-3}y^{3} - \cdots - x^{2}y^{n-2} - xy^{n-1}\Big) 
\end{multline*}
We can see that every term (other than \(x^{n}\) and \(y^{n}\) ) has a corresponding additive inverse, meaning that every such pair
will equal to zero. We are left then with:
\begin{multline*}
	= (x^{n} + x^{n-1}y + x^{n-2}y^2 + x^{n-3}y^3 + \cdots + x^{2}  y^{n-2} + x  y^{n-1})\\
	+ \Big(-y^{n} - x^{n-1}y - x^{n-2}y^{2} - x^{n-3}y^{3} - \cdots - x^{2}y^{n-2} - xy^{n-1}\Big) = x^{n} - y^{n}
\end{multline*}
and that is what had to be proved.

\subsection*{(\rom{6})}

Let us observe that:
\begin{equation*}
	y^{3} = y \cdot y \cdot y = -\Big(-y \cdot (-y \cdot -y)\Big) = -(-y)^{3}
\end{equation*}

We will use this fact and the result from 1.(iv) to prove that \(x^3 + y^3 = (x+y)(x^2 -xy + y^2)\). We can write our formula
in an equivalent way from observation above:
\begin{equation*}
	x^3 + y^3 = x^3 + (-(-y)^{3}) = x^3 - (-y)^{3}
\end{equation*}

Now, we proceed further using \((x-y)(x^2 + xy + y^2) = x^3 - y^3\):
\begin{equation*}
	x^3 + y^3 = x^3 + (-(-y)^{3}) = x^3 - (-y)^{3} = \Big(x-(-y)\Big)(x^2 + x \cdot (-y) + (-y)^2) = (x+y)(x^{2} - xy + y^{2})
\end{equation*}
and that concludes our end result. 


\section*{2}

One may ask why does the below seem to work:

\begin{align*}
	x^{2} &=  xy\\
	x^{2} - y^{2} &=  xy - y^{2}\\
	(x+y)(x-y) &=  y(x - y)\\
	x + y &=  y\\
	2y &=  y\\
	2 &=  1
\end{align*}

Our assumption is that \(x = y\). But this means that there is an illegal operation from equation three to equation four. Namely, 
we try to use \emph{existence of multiplicative inverses} for this operation:

\begin{equation*}
	(x-y) \cdot (x-y)^{-1} = 1
\end{equation*}

but this is not legal, because from our assumption we can see that:

\begin{align*}
	x &= y \\
	x + (-y) &= y + (-y) = 0
\end{align*}
where we used \emph{existence of additive inverses}. This shows that \((x-y)\) does not have the multiplicative inverse because there is
no number \(0^{-1}\) satisfying \(0 \cdot 0^{-1} = 1\)
 
\section*{3}

\subsection*{(i)}

\begin{align*}
	\frac{a}{b} &= ab^{-1} = ab^{-1} \cdot 1 = ab^{-1} \cdot c \cdot c^{-1} = (ac)(b^{-1}c^{-1}) = (ac)(bc)^{-1} = \frac{ac}{bc}
\end{align*}

where, going from the left, we used the definition of the fraction symbol, \emph{existence of multiplicative identity},
\emph{existence of multiplicative inverses}, \emph{associative and commutative laws for multiplication} and 
the result from exercise 1.(iii) (i.e \((ab)^{-1} = a^{-1}b^{-1}\))

\subsection*{(ii)}

We will start with the right side and arrive to the term on the left side of equality:

\begin{equation*}
	\frac{ad + bc}{bd} =
\end{equation*}
\begin{equation*}
	= (ad + bc)(bd)^{-1} \ \ \ \emph{Definition of fraction symbol}
\end{equation*}
\begin{equation*}
	= \Big((ad)(bd)^{-1}\Big) \cdot \Big((bc)(bd)^{-1}\Big) \ \ \ \mbox{\emph{Distributive law}}
\end{equation*}
\begin{equation*}
	= \Big((ad)(b^{-1}d^{-1})\Big) \cdot \Big((bc)(b^{-1}d^{-1})\Big) \ \ \ \mbox{\emph{Equivalent form of inverse}}
\end{equation*}
\begin{equation*}
	= \Big((ad)(d^{-1}b^{-1})\Big) \cdot \Big((cb)(b^{-1}d^{-1})\Big) \ \ \ \mbox{\emph{Commutative law for multiplication}}
\end{equation*}
\begin{equation*}
	= \Big(a((dd^{-1})b^{-1})\Big) \cdot \Big(c((bb^{-1})d^{-1})\Big) \ \ \ \mbox{\emph{Associative law for multiplication}}
\end{equation*}
\begin{equation*}
	= ab^{-1} + cd^{-1} \ \ \ \mbox{\emph{Existence of multiplicative inverses}}
\end{equation*}
\begin{equation*}
	= \frac{a}{b} + \frac{c}{d} \ \ \ \mbox{\emph{Definition of fraction symbol}}
\end{equation*}
Quod erat demonstrandum.


\subsection*{(iii)}

We know that \((ab)\) is invertible because we start with the assumption that \(a \neq 0\) and \(b \neq 0\).
Which means that \(a \cdot b \neq 0\). Hence, \(ab\) is some real number, and from axioms we know that
every real number other than zero has an inverse.  

Let us first use \emph{Existence of multiplicative inverses} to settle the ground for our proof:
\begin{equation*}
	(ab)^{-1}(ab) = 1
\end{equation*}
Now we will use the chain of properties of real numbers:
\begin{equation*}
	(ab)^{-1}(ab) = 1 \cdot 1 \ \ \ \emph{Existence of the multiplicative identity}
\end{equation*}
\begin{equation*}
	(ab)^{-1}(ab) = (a^{-1}a) \cdot (bb^{-1}) \ \ \ \emph{Existence of multiplicative inverses}
\end{equation*}
\begin{equation*}
	(ab)^{-1}(ab) = (a^{-1}b^{-1}) \cdot (ab) \ \ \ \emph{Associative and commutative laws for multiplication (lots of)} 
\end{equation*}
\begin{equation*}
	(ab)^{-1} = (a^{-1}b^{-1}) \ \ \ \emph{Existence of multiplicative inverses and associative law for multiplication} 
\end{equation*}
and our proof is done.

\subsection*{(iv)}

\begin{equation*}
	\frac{a}{b} \cdot \frac{c}{d} 
\end{equation*}
\begin{equation*}
	= (ab^{-1}) \cdot (cd^{-1}) \ \ \ \emph{Definition of fraction symbol}
\end{equation*}
\begin{equation*}
	= (a \cdot (b^{-1} \cdot (cd^{-1}))) \ \ \ \mbox{\emph{Associative law for multiplication}}
\end{equation*}
\begin{equation*}
	= a \cdot ((b^{-1} \cdot c) d^{-1}) \ \ \ \mbox{\emph{Associative law for multiplication}}
\end{equation*}
\begin{equation*}
	= a \cdot ((cb^{-1}))d^{-1}) \ \ \ \mbox{\emph{Commutative law for multiplication}}
\end{equation*}
\begin{equation*}
	= (a(cb^{-1})) \cdot d^{-1} \ \ \ \mbox{\emph{Associative law for multiplication}}
\end{equation*}
\begin{equation*}
	= ((ac)b^{-1}) \cdot d^{-1} \ \ \ \mbox{\emph{Associative law for multiplication}}
\end{equation*}
\begin{equation*}
	= (ac)(b^{-1}d^{-1}) \ \ \ \mbox{\emph{Associative law for multiplication}}
\end{equation*}
\begin{equation*}
	= (ac)(bd)^{-1} \ \ \ \mbox{\emph{Equivalent form of inverse}}
\end{equation*}
\begin{equation*}
	= \frac{ac}{bd} \ \ \ \emph{Definition of fraction symbol}
\end{equation*}

\subsection*{(v)}

We will first prove a helpful observation:

\begin{lemma}
	For any \(d \in \mathbb{R}\) we have:
	\begin{equation*}
		(d^{-1})^{-1} = d
	\end{equation*}
	provided that \(d \neq 0\)
\end{lemma}
\begin{proof}
	Let us first make an observation that the term \((d^{-1})^{-1}\) is \emph{well defined}.
	Suppose for contradiction that \(d^{-1} = 0\). Then using existence of multiplicative inverse and
	existence of multiplicative identity we get:
	\begin{equation*}
		1 = d^{-1}d = 0 \cdot d = 0
	\end{equation*}
	but \(1 \neq 0\), contradiction.

	We shall use now the properties of the reals, setting up the equation:
	\begin{equation*}
		(d^{-1})^{-1}(d^{-1}) = 1 \ \ \ \emph{Existence of multiplicative inverses}
	\end{equation*}
	\begin{equation*}
		\Big((d^{-1})^{-1}(d^{-1})\Big) \cdot d = 1 \cdot d \ \ \ \emph{Multiplication by d}
	\end{equation*}
	\begin{equation*}
		(d^{-1})^{-1}\Big((d^{-1}) \cdot d\Big) = 1 \cdot d \ \ \ \emph{Associative law for multiplication}
	\end{equation*}
	\begin{equation*}
		(d^{-1})^{-1} \cdot 1 = 1 \cdot d \ \ \ \emph{Existence of multiplicative inverses}
	\end{equation*}
	\begin{equation*}
		(d^{-1})^{-1} = d \ \ \ \emph{Existence of multiplicative identity}
	\end{equation*}
\end{proof}

Consider now the main part, we have for \(b, c, d \neq 0\):

\begin{equation*}
	\frac{a}{b} \mathbf{\Bigg/} \frac{c}{d} = (ab^{-1}) \mathbf{\Bigg/} (cd^{-1}) \ \ \ \emph{Definition of the fraction symbol}
\end{equation*}
\begin{equation*}
	= (ab^{-1}) \cdot (cd^{-1})^{-1} \ \ \ \emph{Definition of the fraction symbol}
\end{equation*}
\begin{equation*}
	= (ab^{-1}) \cdot (c^{-1}  (d^{-1})^{-1}) \ \ \ \emph{Property exercise 1.3.IV}
\end{equation*}
Using now lemma above:
\begin{equation*}
	= (ab^{-1}) \cdot (c^{-1}d)
\end{equation*}
\begin{equation*}
	= (ab^{-1}) \cdot (dc^{-1}) \ \ \ \emph{Commutative law for multiplication}
\end{equation*}
\begin{equation*}
	= (a(b^{-1}d)) \cdot c^{-1} \ \ \ \emph{Associative law for multiplication}
\end{equation*}
\begin{equation*}
	= (a(db^{-1})) \cdot c^{-1} \ \ \ \emph{Commutative law for multiplication}
\end{equation*}
\begin{equation*}
	= ((ad)b^{-1}) \cdot c^{-1} \ \ \ \emph{Associative law for multiplication}
\end{equation*}
\begin{equation*}
	= (ad) \cdot (b^{-1}c^{-1}) \ \ \ \emph{Associative law for multiplication}
\end{equation*}
\begin{equation*}
	= (ad) \cdot (bc)^{-1} \ \ \ \emph{Property exercise 1.3.IV}
\end{equation*}
\begin{equation*}
	= \frac{ad}{bc} \ \ \ \emph{Definition of the fraction symbol}
\end{equation*}

\subsection*{(vi)}

Assume \(ad = bc\) with \(b, d \neq 0\). We will use the properties of the reals:

\begin{equation*}
	(ad) \cdot b^{-1} = (bc) \cdot b^{-1} \ \ \ \emph{Multiplying by} \ b^{-1}
\end{equation*}
\begin{equation*}
	\Big((ad) \cdot b^{-1}\Big) \cdot d^{-1} = \Big((bc) \cdot b^{-1}\Big) \cdot d^{-1} \ \ \ \emph{Multiplying by} \ d^{-1}
\end{equation*}
\begin{equation*}
	(ab^{-1}) \cdot dd^{-1} = (cd^{-1}) \cdot bb^{-1} \ \ \ \emph{Assoc. and commut. law for multiplication}
\end{equation*}
\begin{equation*}
	(ab^{-1}) \cdot 1 = (cd^{-1}) \cdot 1 \ \ \ \emph{Existence of multiplicative inverses}
\end{equation*}
\begin{equation*}
	ab^{-1}  = cd^{-1} \ \ \ \emph{Existence of a multiplicative identity}
\end{equation*}
\begin{equation*}
	\frac{a}{b} = \frac{c}{d} \ \ \ \emph{Definition of fraction symbol}
\end{equation*}

For the other direction, assume \(\frac{a}{b} = \frac{c}{d}\) with \(b, d \neq 0\). We have:

\begin{equation*}
	ab^{-1} = cd^{-1} \ \ \ \emph{Definition of fraction symbol}
\end{equation*}
\begin{equation*}
	(ab^{-1}) \cdot d = (cd^{-1}) \cdot d \ \ \ \emph{Multiplying by d}
\end{equation*}
\begin{equation*}
	\Big((ab^{-1}) \cdot d\Big) \cdot b = \Big((cd^{-1}) \cdot d\Big) \cdot b \ \ \ \emph{Multiplying by b}
\end{equation*}
\begin{equation*}
	(ad) \cdot (b^{-1}b) = (cb) \cdot (dd^{-1}) \ \ \ \emph{Assoc. and commut. law for multiplication}
\end{equation*}
\begin{equation*}
	(ad) \cdot 1 = (cb) \cdot 1 \ \ \ \emph{Existence of multiplicative inverses}
\end{equation*}
\begin{equation*}
	ad = cb  \ \ \emph{Existence of a multiplicative identity}
\end{equation*}
\begin{equation*}
	ad = bc  \ \ \emph{Commutative law for multiplication}
\end{equation*}

Assume now \(\frac{a}{b} = \frac{b}{a}\). Let us write this equation in an equivalent way:

\begin{equation*}
	aa = bb
\end{equation*}
\begin{equation*}
	a^{2} = b^{2}
\end{equation*}
\begin{equation*}
	a^{2} - b^{2} = (a-b)(a+b) = 0
\end{equation*}

So either \(a - b = 0\) or \(a + b = 0\). Thus \(\frac{a}{b} = \frac{b}{a}\) if:

\begin{equation*}
	a = b \ \ \mbox{or} \ \ a = -b
\end{equation*}

\section*{5}

\subsection*{(i)}

Assume \(a < b\) and \(c < d\). Observe that \((b - a)\) and \((d - c)\) are
positive numbers. The positive numbers are closed under addition so \((b - a) + (d - c)\)
is also a positive number. By definition we have then:
\begin{equation*}
	(b - a) + (d - c) > 0
\end{equation*} 
\begin{equation*}
	\Big((b - a) + (d - c)\Big) + a > 0 + a
\end{equation*}
\begin{equation*}
	\Bigg(\Big((b - a) + (d - c)\Big) + a\Bigg) + c > (0 + a) + c
\end{equation*}

It follows from Associative and commutative law for addition and existence of additive inverses and an
additive identity that:

\begin{equation*}
	b + d > a + c
\end{equation*}


\subsection*{(ii)}

Suppose \(a < b\). Add \(-a\) to both sides:

\begin{equation*}
	a + (-a) < b + (-a)
\end{equation*}
by existence of additive inverses:
\begin{equation*}
	0 < b + (-a) = b - a
\end{equation*}
Adding \(-b\) to both sides yields:
\begin{equation*}
	-b < (b - a) + (-b)
\end{equation*}
And again applying existence of additive inverses on top off associative and commutative law for addition we finally get:
\begin{equation*}
	-b < -a
\end{equation*}

\subsection*{(iii)}

Suppose \(a < b\) and \(c > d\). Observe that \(b - a\) and \(c - d\) are positive numbers. Since the positive numbers are closed
under addition it follows that:

\begin{equation*}
	(b - a) + (c - d) > 0
\end{equation*}

Adding (a - c) to both sides yields:

\begin{equation*}
	\Big((b - a) + (c - d)\Big) + (a - c) > a - c
\end{equation*}

It follows from Associative and commutative law for addition and existence of additive inverses and an
additive identity that:

\begin{equation*}
	b - d > a - c
\end{equation*}

\subsection*{(iv)}

Suppose \(a < b\) and \(c > 0\). Observe that \(b - a\) is a positive number.
Since the positive numbers are closed under multiplication it follows that:

\begin{equation*}
	(b-a) \cdot c > 0
\end{equation*}
Applying distributive law we get:

\begin{equation*}
	bc + (-ac) > 0
\end{equation*}

and adding \(ac\) to both sides we end up with:

\begin{equation*}
	bc + (-ac) + (ac) > 0 + ac 
\end{equation*}

It follows from existence of additive inverses and existence of an additive identity that:

\begin{equation*}
	bc > ac
\end{equation*}

\subsection*{(v)}

Suppose \(a < b\) and \(c < 0\). Observe that \(b - a\) and \(-c\) are positive numbers.
Since the positive numbers are closed under multiplication it follows that:

\begin{equation*}
	(b-a) \cdot -c > 0
\end{equation*}
Applying distributive law we get:

\begin{equation*}
	-bc + (-a)(-c) = -bc + ac > 0
\end{equation*}

and adding \(bc\) to both sides we end up with:

\begin{equation*}
	-bc + ac + bc > 0 + bc
\end{equation*}

It follows from existence of additive inverses and existence of an additive identity that:

\begin{equation*}
	ac > bc
\end{equation*}

\subsection*{(vi)}

Suppose \(a > 1\). Observe that \(a\) and \(a - 1\) are positive numbers.
Since the positive numbers are closed under multiplication it follows that:

\begin{equation*}
	(a -1) \cdot a > 0
\end{equation*}

Applying distributive law we get:

\begin{equation*}
	a^2 + (-a) = a^{2} - a > 0
\end{equation*}

Adding to both sides \(a\) it follows from existence of additive inverses and existence of an additive identity that:

\begin{equation*}
	a^2 > a
\end{equation*}

\subsection*{(vii)}

Suppose \(0 < a < 1\). Observe that
\((1 - a)\) is a positive number. Since \(a\) is also a positive number we use closure
of positive numbers under multiplication to get:

\begin{equation*}
	(1 -a) \cdot a > 0
\end{equation*}

Using distributive law we get:

\begin{equation*}
	a + (-a^{2}) > 0
\end{equation*}

And from here we conclude:

\begin{equation*}
	a > a^{2}
\end{equation*}

\subsection*{(viii)}

Suppose \(0 \leq a < b\) and \(0 \leq c < d\). Using the result from exercise 5.IV we
can write the assumption in an equivalent way:

\begin{equation*}
	ac \leq bc \ \ \mbox{and} \ \ bc < bd
\end{equation*}

Now, observe that \(bc - ac\) can be positive or zero but \(bd - bc\) is a positive
number. Adding those two together will yield a positive number, so:

\begin{equation*}
	(bd - bc) + (bc - ac) > 0
\end{equation*}

but this is just:

\begin{equation*}
	bd - ac > 0
\end{equation*}

and we conclude with:

\begin{equation*}
	bd > ac
\end{equation*}

\subsection*{(IX)}

Suppose \(0 \leq a < b\). Using the result from exercise 5.VIII we immediately have:

\begin{equation*}
	a \cdot a = a^{2} < b \cdot b = b^{2}
\end{equation*}

\subsection*{(X)}

Suppose \(a, b \geq 0\) and \(a^{2} < b^{2}\). Assume for contradiction that \(a \geq b\). From the result of 5.IX we
get that:

\begin{equation*}
	a \geq b \implies a^{2} \geq b^{2}
\end{equation*}

But clearly \(a^{2} \geq b^{2}\) and \(a^{2} < b^{2}\) cannot coexist together.

\section*{12}

\subsection*{(I)}

We will prove the statement by considering cases where \(x, y\) can have different or the same signs.

Consider \(x \geq 0\) and \(y \geq 0\). Notice that \(xy \geq 0\) We have:

\begin{equation*}
	|xy| = x \cdot y = |x| \cdot |y|
\end{equation*}

Suppose now that \(x < 0\) and \(y < 0\). Notice that \(xy > 0\) We have:

\begin{equation*}
	|xy| = x \cdot y = (-x) \cdot (-y) = |x||y|
\end{equation*}

For the last case, assume without loss of generality that \(x \geq 0\) and \(y < 0\). If \(x = 0\) we
have \(|0| = 0 = 0 \cdot 0 = |0||0|\). If \(x > 0\), then:

\begin{equation*}
	|xy| = -(xy) = -(yx) = (-y)x = |y||x| = |x||y|
\end{equation*}

\subsection*{(II)}

Assume \(x \neq 0\). We have:

\begin{equation*}
	1 = |1| = |xx^{-1}| = |x||x^{-1}| = |x|\abs{\frac{1}{x}}
\end{equation*}
It follows from:
\begin{equation*}
	1 \cdot |x|^{-1} = \Big(|x|\abs{\frac{1}{x}}\Big)|x|^{-1} = \Big(|x||x|^{-1}\Big)\abs{\frac{1}{x}} = \abs{\frac{1}{x}}
\end{equation*}
\begin{equation*}
	|x|^{-1} = \frac{1}{|x|} =  \abs{\frac{1}{x}}
\end{equation*}

\subsection*{(III)}

Assume \(y \neq 0\). We have:

\begin{equation*}
	\abs{\frac{x}{y}} = |xy^{-1}| = |x||y^{-1}| = |x|\abs{\frac{1}{y}}
\end{equation*}
Using now the previous result (12.(II)):

\begin{equation*}
	= |x|\frac{1}{\abs{y}} = |x||y|^{-1} = \frac{\abs{x}}{\abs{y}}
\end{equation*}

and that completes the proof.

\subsection*{(IV)}

Suppose \(x, y\in \mathbb{R}\). We have:

\begin{equation*}
	|x - y| = |(x + 0) + (-y + 0)| \ \ \ \emph{Existence of an additive identity}
\end{equation*}
\begin{equation*}
	\leq |(x + 0)| + |(-y + 0)| \ \ \ \emph{Triangle inequality}
\end{equation*}
\begin{equation*}
	= |x| + |-y| \ \ \ \emph{Existence of an additive identity}
\end{equation*}
\begin{equation*}
	= |x| + |(-1) \cdot y| \ \ \ \emph{Using a fact that (-a)b = -(ab)}
\end{equation*}
\begin{equation*}
	= |x| + |-1||y| \ \ \ \emph{Multiplication of absolute values}
\end{equation*}
\begin{equation*}
	= |x| + 1|y| \ \ \ \emph{Absolute value definition}
\end{equation*}
\begin{equation*}
	= |x| + |y| \ \ \ \emph{Existence of a multiplicative identity}
\end{equation*}

\subsection*{(V)}

Suppose \(x, y \in \mathbb{R}\). We have:

\begin{equation*}
	|x| = |x + 0| \ \ \ \emph{Existence of an additive identity}
\end{equation*}
\begin{equation*}
	= |x + ((-y) + y)| \ \ \ \emph{Existence of additive inverses}
\end{equation*}
\begin{equation*}
	= |(x + (-y)) + y| \ \ \ \emph{Associative law for addition}
\end{equation*}
\begin{equation*}
	\leq |x + (-y)| + |y| \ \ \ \emph{Triangle inequality}
\end{equation*}
So:
\begin{equation*}
	|x| \leq |x - y| + |y|
\end{equation*}

And adding \(-|y|\) and using existence of additive inverses and existence of an additive identity:

\begin{equation*}
	|x| - |y| \leq |x - y| 
\end{equation*}

and that completes the proof.

\subsection*{(VI)}

\begin{lemma}[\textbf{1}]
	Suppose \(a, b \in \mathbb{R}\). We have the following fact:
	\begin{equation*}
		-b \leq a \leq b \iff |a| \leq b
	\end{equation*}
	% Moreover: \(-|a| \leq a \leq |a|\)
\end{lemma}
\begin{proof}
	Suppose \(a \geq 0\) and \(-b \leq a \leq b\). From \(|a| = a\) clearly we have \(-b \leq |a| \leq b\).
	For \(a < 0\), observe:
	\begin{equation*}
		-b \leq a \implies -a = |a| \leq b
	\end{equation*}
	In other direction, assume \(|a| \leq b\). For the case \(a \geq 0\) observe that \(b \geq 0 \geq -b\) so that:
	\begin{equation*}
		-b \leq 0 \leq |a| = a \leq b
	\end{equation*}
	For \(a < 0\) we have:
	\begin{equation*}
		0 < |a| = -a \leq b
	\end{equation*}
	which means that \(a \geq -b\). Observe that both \(-a\) and \(b\) are positive numbers. Then \(b - a > 0\) and it follows:
	\begin{equation*}
		b > a \geq -b
	\end{equation*}
	And since strict inequalities imply weak we conclude with:
	\begin{equation*}
		-b \leq a \leq b
	\end{equation*}
\end{proof}

\begin{lemma}[\textbf{2}]
	For every \(x \in \mathbb{R}\) we have \(|x| = |-x|\)
\end{lemma}
\begin{proof}
	Using elementary operations on the absolute value:
	\begin{equation*}
		|-x| = |(-1) \cdot x| = |-1||x| = 1|x| = |x|
	\end{equation*}
\end{proof}
\newpage

Suppose \(x, y \in \mathbb{R}\). Let us start with the property \(12.(V)\):

\begin{equation*}
	|x| - |y| \leq |x-y|
\end{equation*}

but also:

\begin{equation*}
	-|x - y| \leq |y| - |x|
\end{equation*}

Again, applying the property \(12.(V)\):

\begin{equation*}
	|y| - |x| \leq |y - x|
\end{equation*}

Using Lemma(2), we observe:

\begin{equation*}
	|y - x| = |-(x - y)| = |x-y|
\end{equation*}

Then combining all we have:

\begin{equation*}
	-|x - y| \leq |y| - |x| \leq |y - x| = |x-y|
\end{equation*}

Now using Lemma(1) we quickly notice:

\begin{equation*}
	-|x - y| \leq |y| - |x| \leq |x-y| \iff \abs{|y| - |x|} \leq |x-y|
\end{equation*}

And now Lemma(2) again:

\begin{equation*}
	\abs{|y| - |x|} = \abs{-(|y| - |x|)} = \abs{|x| - |y|}  \leq |x-y|
\end{equation*}

and the proof is complete.

\subsection*{(VII)}

Suppose \(x, y, z \in \mathbb{R}\). Applying the triangle inequality twice:

\begin{equation*}
	|x + y + z| = |(x + y) + z| \leq |x + y| + |z| \leq |x| + |y| + |z|
\end{equation*}
gets us the proposition. The equality occurs when \(x,y,z\) are all non-negative or negative because we have then either:

\begin{equation*}
	|x  + y + z| = x + y + z = |x| + |y| + |z| \ \ \mbox{for \(x, y ,z \geq 0\)}
\end{equation*}
or:
\begin{equation*}
	|x + y + z| = -(x + y + z) = (-x) + (-y) + (-z) = |x| + |y| + |z| \ \ \mbox{for \(x, y ,z < 0\)}
\end{equation*}


\section*{13}

First, observe that \(\max(x, y)\) and \(\min(x, y)\) definitions can be written like this:

\begin{equation*}
	\max(x, y) = 
	\begin{cases}
	x, \ \ \mbox{if} \ \ x \geq y \\
	y, \ \ \mbox{if} \ \ x < y
	\end{cases}
\end{equation*}

and

\begin{equation*}
	\min(x, y) = 
	\begin{cases}
		x, \ \ \mbox{if} \ \ x \leq y \\
		y, \ \ \mbox{if} \ \ x > y
	\end{cases}
\end{equation*}
a quick interaction with those objects confirms that they are well-defined. Now, suppose
\(x, y \in \mathbb{R}\) and suppose further that \(x \geq y\) such that \(\max(x, y) = x\).
We shall notice that \(x - y \geq 0\). With this in mind, we can write that \(x - y = |x - y|\).


\begin{equation*}
	x - y = |x - y|
\end{equation*}
\begin{equation*}
	x = |x - y| + y
\end{equation*}
\begin{equation*}
	2x - x = |x - y| + y
\end{equation*}
\begin{equation*}
	2x = |x - y| + y + x
\end{equation*}
\begin{equation*}
	x = \frac{\abs{x - y} + y + x}{2}
\end{equation*}
Remembering now that \(|x| = |-x|\) we conclude with:
\begin{equation*}
	x = \max(x, y) = \frac{\abs{y - x} + y + x}{2}
\end{equation*}
Consider now \(x < y\) such that \(\max(x, y) = y\). We again notice \(y - x > 0\) which implies
\(y - x = |y - x|\). Tying everything together, we get:

\begin{equation*}
	y - x = |y - x|
\end{equation*}
\begin{equation*}
	y = |y - x| + x
\end{equation*}
\begin{equation*}
	2y - y = |y - x| + x
\end{equation*}
\begin{equation*}
	2y = |y - x| + x + y
\end{equation*}
\begin{equation*}
	\max(x, y) = y = \frac{|y - x| + x + y}{2}
\end{equation*}


A similar argument is applied to \(\min(x, y)\). Suppose \(x, y \in \mathbb{R}\) 
and suppose further that \(x \leq y\) such that \(\min(x, y) = x\). Then \(y - x \geq 0\)
and \(y - x = |y - x|\). It follows:

\begin{equation*}
	y - x = |y - x|
\end{equation*}
\begin{equation*}
	x = y - |y - x|
\end{equation*}
\begin{equation*}
	2x - x = y - |y - x|
\end{equation*}
\begin{equation*}
	2x = y + x - |y - x|
\end{equation*}
\begin{equation*}
	\min(x, y) = x = \frac{y + x - |y - x|}{2}
\end{equation*}

And finally, suppose that \(x > y\) such that \(\min(x, y) = y\). Then:

\begin{equation*}
	|y - x| = |x - y| = x - y
\end{equation*}
\begin{equation*}
	y = x - |y - x|
\end{equation*}
\begin{equation*}
	2y - y = x - |y - x|
\end{equation*}
\begin{equation*}
	2y = x - |y - x| + y
\end{equation*}
\begin{equation*}
	\min(x, y) = y = \frac{x + y - |y - x|}{2}
\end{equation*}

Observe now, that \(\max(x, y, z) = \max(x, \max(y, z))\) is well-defined, simply because we can mirror
the definition used for a two-variable case by putting:

\begin{equation*}
	\max(x, y, z) = 
	\begin{cases}
	x, \ \ \mbox{if} \ \ x \geq \max(y, z) \\
	\max(y, z), \ \ \mbox{if} \ \ x < \max(y, z)
	\end{cases}
\end{equation*}

and

\begin{equation*}
	\min(x, y, z) = 
	\begin{cases}
		x, \ \ \mbox{if} \ \ x \leq \min(y, z) \\
		\min(y, z), \ \ \mbox{if} \ \ x > \min(y, z)
	\end{cases}
\end{equation*}
treating the output of \(\min(y,z)\) and \(\max(y,z)\) as a single variable.
We can then check first for \(\max(x, y, z)\):


\begin{equation*}
	\max(x, y, z) = \max(x, \max(y, z))
\end{equation*}
\begin{equation*}
	= \frac{x + \max(y, z) + \abs{\max(y,z) - x} }{2}
\end{equation*}
\begin{equation*}
	= \frac{x + \frac{y + z + \abs{z - y}}{2} + \abs{\frac{y + z + \abs{z - y}}{2} - x} }{2}
\end{equation*}
\begin{equation*}
	= \frac{\frac{2x + y + z + \abs{z - y}}{2} + \abs{\frac{y + z + \abs{z - y} -2x}{2}} }{2}
\end{equation*}
\begin{equation*}
	= \frac{ \frac{1}{2} \Bigg( \Big(2x + y + z + \abs{z - y}\Big) + \Big(\abs{y + z + \abs{z - y} -2x}\Big)\Bigg)}{2}
\end{equation*}
\begin{equation*}
	= \frac{ \Bigg( \Big(2x + y + z + \abs{z - y}\Big) + \Big(\abs{y + z + \abs{z - y} -2x}\Big)\Bigg)}{4}
\end{equation*}
and for \(\min(x, y, z)\):
\begin{equation*}
	\min(x, y, z) = \min(x, \min(y, z))
\end{equation*}
\begin{equation*}
	= \frac{x + \min(y, z) - \abs{\min(y, z) - x}}{2}
\end{equation*}
\begin{equation*}
	= \frac{x + \frac{y + z - \abs{z - y}}{2} - \abs{\frac{y + z - \abs{z - y}}{2} - x}}{2}
\end{equation*}
\begin{equation*}
	= \frac{\frac{2x + y + z - \abs{z - y}}{2} - \abs{\frac{y + z - \abs{z - y} - 2x}{2}}}{2}
\end{equation*}
\begin{equation*}
	= \frac{\frac{1}{2}\Bigg( \Big(2x + y + z - \abs{z - y}\Big) - \Big(\abs{y + z - \abs{z - y} - 2x}\Big)\Bigg)}{2}  
\end{equation*}
\begin{equation*}
	= \frac{\Big(2x + y + z - \abs{z - y}\Big) - \Big(\abs{y + z - \abs{z - y} - 2x}\Big)}{4}  
\end{equation*}
and that completes the proof.


\section*{14}
\subsection*{(a)}

Assume \(a \in \mathbb{R}\). From the definition of the absolute value:

\begin{equation*}
	|-a| = |(-1)a| = |-1||a| = 1|a| = |a|
\end{equation*}
where we also used the result from (12.(I)).

\subsection*{(b)}

Suppose \(a \geq 0\) and \(-b \leq a \leq b\). From \(|a| = a\) clearly we have \(-b \leq |a| \leq b\).
For \(a < 0\), observe:
\begin{equation*}
	-b \leq a \implies -a = |a| \leq b
\end{equation*}
\newline
In other direction, assume \(|a| \leq b\). For the case \(a \geq 0\) observe that \(b \geq 0 \geq -b\) so that:
\begin{equation*}
	-b \leq 0 \leq |a| = a \leq b
\end{equation*}
For \(a < 0\) we have:
\begin{equation*}
	0 < |a| = -a \leq b
\end{equation*}
which means that \(a \geq -b\). Observe that both \(-a\) and \(b\) are positive numbers. Then \(b - a > 0\) and it follows:
\begin{equation*}
	b > a \geq -b
\end{equation*}
And since strict inequalities imply weak we conclude with:
\begin{equation*}
	-b \leq a \leq b
\end{equation*}
Suppose now that \(a \in \mathbb{R}\). Then trivially we have:
\begin{equation*}
	|a| \leq |a|
\end{equation*}
as this relationship is always equal. Applying the statement from above:
\begin{equation*}
	-|a| \leq a \leq |a|
\end{equation*}
completes the proof.

\subsection*{(c)}
Suppose that \(a, b \in \mathbb{R}\). Let us set up inequalities for \(a\) and \(b\) using the proven fact before:

\begin{equation*}
	-|a| \leq a \leq |a|
\end{equation*}
and
\begin{equation*}
	-|b| \leq b \leq |b|
\end{equation*}
By adding those two together we end up with:
\begin{equation*}
	-(|a| + |b|) \leq a + b \leq (|a| + |b|)
\end{equation*}
And applying again the property from (b) we conclude with:
\begin{equation*}
	|a + b| \leq |a| + |b|
\end{equation*}

\section*{18}

\subsection*{(a)}

Suppose that \(b^{2} - 4c \geq 0\). We shall see that both:

\begin{equation*}
	\frac{-b + \sqrt{b^{2} - 4c}}{2} \ \ \mbox{and} \ \ \frac{-b - \sqrt{b^{2} - 4c}}{2}
\end{equation*}

satisfy the equation \(x^{2} + bx + c = 0\). Plugging in the first formula we have:

\begin{equation*}
	x^{2} + bx + c = \Bigg(\frac{-b + \sqrt{b^{2} - 4c}}{2}\Bigg)^{2} + b\Bigg(\frac{-b + \sqrt{b^{2} - 4c}}{2}\Bigg) + c
\end{equation*}
\begin{equation*}
	= \Bigg(\frac{(b^{2} - 4c) - 2\Big(\sqrt{b^{2} - 4c}\Big)b + b^{2}}{4}\Bigg) + b\Bigg(\frac{-2b + 2\sqrt{b^{2} - 4c}}{4}\Bigg) + c
\end{equation*}
\begin{equation*}
	= \Bigg(\frac{(b^{2} - 4c) - 2\Big(\sqrt{b^{2} - 4c}\Big)b + b^{2}}{4}\Bigg) + \Bigg(\frac{-2b^{2} + 2b\sqrt{b^{2} - 4c}}{4}\Bigg) + c
\end{equation*}
\begin{equation*}
	= \Bigg(\frac{(b^{2} - 4c) - 2\Big(\sqrt{b^{2} - 4c}\Big)b + b^{2}}{4}\Bigg) + \Bigg(\frac{-2b^{2} + 2b\sqrt{b^{2} - 4c}}{4}\Bigg) + \frac{4c}{4}
\end{equation*}
\begin{equation*}
	= \frac{(b^{2} - 4c) - 2\Big(\sqrt{b^{2} - 4c}\Big)b + b^{2} - 2b^{2} + 2b\sqrt{b^{2} - 4c} + 4c}{4}
\end{equation*}
\begin{equation*}
	= \frac{b^{2} - 4c - 2b\sqrt{b^{2} - 4c} + b^{2} - 2b^{2} + 2b\sqrt{b^{2} - 4c} + 4c}{4}
\end{equation*}
\begin{equation*}
	= \frac{b^{2} + b^{2} - 2b^{2} - 4c + 4c - 2b\sqrt{b^{2} - 4c} + 2b\sqrt{b^{2} - 4c}}{4}
\end{equation*}
\begin{equation*}
	= \frac{(b^{2} + b^{2} - 2b^{2}) + (- 4c + 4c) + \Big(-2b\sqrt{b^{2} - 4c} + 2b\sqrt{b^{2} - 4c}\Big)}{4} = 0
\end{equation*}

And for the second formula we also have:

\begin{equation*}
	x^{2} + bx + c = \Bigg(\frac{-b - \sqrt{b^{2} - 4c}}{2}\Bigg)^{2} + b\Bigg(\frac{-b - \sqrt{b^{2} - 4c}}{2}\Bigg) + c
\end{equation*}
\begin{equation*}
	= \Bigg(\frac{b + \sqrt{b^{2} - 4c}}{2}\Bigg)^{2} + b\Bigg(\frac{-b - \sqrt{b^{2} - 4c}}{2}\Bigg) + c
\end{equation*}
\begin{equation*}
	= \frac{b^{2} + 2b\sqrt{b^{2} - 4c} + (b^{2} - 4c)}{4} + \frac{-2b^{2} - 2b\sqrt{b^{2} - 4c}}{4} + \frac{4c}{4}
\end{equation*}
\begin{equation*}
	= \frac{b^{2} + 2b\sqrt{b^{2} - 4c} + (b^{2} - 4c) -2b^{2} - 2b\sqrt{b^{2} - 4c} + 4c}{4}
\end{equation*}
\begin{equation*}
	= \frac{(b^{2} + b^{2} - 2b^{2}) + (2b\sqrt{b^{2} - 4c} - 2b\sqrt{b^{2} - 4c}) + (-4c + 4c)}{4} = 0
\end{equation*}
and we are done.


\subsection*{(b)}

Suppose \(b^{2} - 4c < 0\). Notice that \(\frac{b^{2}}{4} < c\). Recall also that for any \(a \in \mathbb{R}\), \(a^{2} \geq 0\). We have:

\begin{equation*}
	x^{2} + bx + c > x^{2} + bx + \frac{b^{2}}{4} = (x + \frac{b}{2})^{2} \geq 0
\end{equation*}
and thus no \(x\) obeys the equation \(x^{2} + bx + c = 0\).

\subsection*{(c)}

Assume first that \(x = 0\) and \(y \neq 0\). We immediately get:
\begin{equation*}
	x^{2} + xy + y^{2} = 0 + 0 \cdot y + y^{2} = y^{2} > 0 
\end{equation*}
which holds always for every \(y \in \mathbb{R} \setminus 0\).

Similarly, assume \(x \neq 0\) and \(y = 0\). Then:
\begin{equation*}
	x^{2} + xy + y^{2} = x^{2} + x \cdot 0 + 0^{2} = x^{2} > 0
\end{equation*}
For the most interesting case assume that \(x \neq 0\) and \(y \neq 0\). Observe that:

\begin{equation*}
	y^{2} - 4y^{2} = y^{2}(1 - 4) = (-3)y^{2} < 0
\end{equation*}

Because \(y^{2} - 4y^{2}\) is not a positive number, one can use the result from (18.b), namely,
let \(b = y\) and \(c = y^{2}\). Now \(b^{2} - 4c < 0\). We have 
\begin{equation*}
	x^{2} + xy + y^{2} = x^{2} + xb + c > 0
\end{equation*}
And this finishes the proof.

\end{document}
